%% Question bank chapter file - generated by bank_tools.py build (do not edit while the bank is still being built:
%% it is regenerated from the per-test data in Question Bank\data). Every question is one bankq block:
%%   \begin{bankq}[id=..., source={...}, paper=.., qno=.., type=..., subtopic={...}, answer={...}, answer_src=key|solved,
%%                 working={...}, difficulty=1-10, diff_parts={K S I C F}, diff_note={...},
%%                 group=<passage id or empty>, also_in={duplicate copies in other tests}]
%% difficulty = min(10, K+S+I+C+F-1): K concepts 1-3, S steps 1-4, I insight 0-2, C computation 0-2, F format 0-1
%%   ...question LaTeX, options, figures (\bfig / \bfigi -> img/)...
%%   \end{bankq}
%% Passages shared by several questions are bankpassage blocks; their questions carry group=<passage id>.
\documentclass[10pt]{article}
\usepackage[a4paper,margin=1.8cm]{geometry}
\usepackage[T1]{fontenc}
\usepackage{amsmath,amssymb,mathtools}
\usepackage{graphicx}
\usepackage[export]{adjustbox}
\usepackage{xcolor,needspace,keyval}
\usepackage{enumitem}
\usepackage{booktabs,array,tabularx}
\usepackage[hidelinks]{hyperref}
\setlength{\parindent}{0pt}
\setlength{\parskip}{3pt}
\makeatletter
\define@key{bq}{id}{\def\bq@id{#1}}\define@key{bq}{source}{\def\bq@src{#1}}\define@key{bq}{paper}{\def\bq@paper{#1}}
\define@key{bq}{qno}{\def\bq@qno{#1}}\define@key{bq}{type}{\def\bq@type{#1}}\define@key{bq}{subtopic}{\def\bq@sub{#1}}
\define@key{bq}{answer}{\def\bq@ans{#1}}\define@key{bq}{answer_src}{\def\bq@asrc{#1}}\define@key{bq}{working}{\def\bq@work{#1}}
\define@key{bq}{group}{\def\bq@group{#1}}\define@key{bq}{also_in}{\def\bq@also{#1}}
\define@key{bq}{difficulty}{\def\bq@diff{#1}}\define@key{bq}{diff_parts}{\def\bq@dparts{#1}}\define@key{bq}{diff_note}{\def\bq@dnote{#1}}
\newenvironment{bankq}[1][]{\def\bq@work{}\def\bq@also{}\def\bq@group{}\def\bq@diff{?}\def\bq@dparts{}\def\bq@dnote{}\setkeys{bq}{#1}%
  \par\needspace{4\baselineskip}\medskip\noindent{\small\textbf{\bq@id}\hfill\color{gray}[\bq@src, Paper \bq@paper, Q\bq@qno]
  \bq@type\ $\cdot$ \bq@sub\ $\cdot$ \textbf{D\bq@diff}}\par\nopagebreak\relax}
  {\par{\small\color{teal}Answer: \textbf{\bq@ans}\ifx\bq@work\@empty\else\ (solved: \bq@work)\fi
  \ifx\bq@also\@empty\else\par\color{gray}Also in: \bq@also\fi}\par}
\define@key{bp}{id}{\def\bp@id{#1}}\define@key{bp}{source}{\def\bp@src{#1}}\define@key{bp}{questions}{\def\bp@qs{#1}}
\newenvironment{bankpassage}[1][]{\setkeys{bp}{#1}\par\needspace{4\baselineskip}\bigskip\noindent
  {\small\textbf{Passage \bp@id}\hfill\color{gray}[\bp@src; questions \bp@qs]}\par\nopagebreak\relax\itshape}{\par}
\makeatother
\newcommand{\bfig}[2]{\par\smallskip{\centering\includegraphics[width=#1pt,max width=0.62\linewidth,max height=0.32\textheight]{img/#2}\par}\smallskip}
\newcommand{\bfigi}[2]{\includegraphics[width=#1pt,max width=\linewidth,max height=0.32\textheight,valign=c]{img/#2}}
\newcommand{\bmissing}[1]{\textit{[figure missing in the source conversion: \detokenize{#1}]}}
\newenvironment{qoptions}{\begin{itemize}[leftmargin=1.8em,itemsep=1pt,topsep=2pt,label={}]}{\end{itemize}}
\newcommand{\opt}[2]{\item[(#1)] #2}
\providecommand{\fwstart}{}\providecommand{\fwend}{}\providecommand{\qfullsplit}{}
\newcolumntype{L}[1]{>{\raggedright\arraybackslash}p{#1}}
\newcolumntype{C}[1]{>{\centering\arraybackslash}p{#1}}
\newcolumntype{R}[1]{>{\raggedleft\arraybackslash}p{#1}}
\providecommand{\sectionhdr}[1]{}\providecommand{\subsectionhdr}[1]{}


\begin{document}
{\Large\bfseries Amines --- Question Bank}\par\smallskip
{\small Chemistry chapter C22. Questions from all converted test papers, grouped by source series. Answers: printed key where it matches the paper, otherwise solved (answer\_src=solved, with working). Exact repeats across papers appear once, with also\_in listing the other copies.}\par\medskip\hrule

\section*{Score Advanced papers}
\begin{bankq}[id=SA01-P2-CHEM-S1-Q03, source={Score Advanced - Test 01 (22 Sep 2020)}, paper=2, qno=3, type=multi, subtopic={Tests}, answer={A,B,C,D}, answer_src=key, difficulty=7, diff_parts={K2 S2 I2 C1 F1}, diff_note={concepts: reaction sequence with amines, tests | steps: intermediates; carbylamine; products}]
Consider the reaction given below :\par
\bfig{317.0}{SA01_p2_chem_s1_q03_1.png}

Select the correct statement(s)
\begin{qoptions}
\opt{A}{(x) is \bfigi{70.9}{SA01_p2_chem_s1_q03_2.png}}
\opt{B}{The product (y) is gives carbayl amine test}
\opt{C}{(z) is \bfigi{70.9}{SA01_p2_chem_s1_q03_3.png}}
\opt{D}{conversion of (x) into (y) involves nucleophilic substitution}
\end{qoptions}
\end{bankq}
\begin{bankq}[id=SA01-P2-CHEM-S1-Q07, source={Score Advanced - Test 01 (22 Sep 2020)}, paper=2, qno=7, type=multi, subtopic={Basicity}, answer={B,C}, answer_src=key, difficulty=3, diff_parts={K1 S1 I1 C0 F1}, diff_note={concepts: dealkylation of tertiary amines | steps: von Braun type reaction; reagent selection}]
Which of the following reagent(s) would be most suitable to convert $\left(\mathrm{CH_3}\right)_3\mathrm N$ into $\left(\mathrm{CH_3}\right)_2\mathrm{NH}$
\begin{qoptions}
\opt{A}{HCl only}
\opt{B}{NCBr}
\opt{C}{H$_3$O$^+$;$\Delta$}
\opt{D}{(PPh$_3$ + H$_2$O) only}
\end{qoptions}
\end{bankq}
\begin{bankq}[id=SA02-P2-CHEM-S1-Q05, source={Score Advanced - Test 02 (13 May 2021)}, paper=2, qno=5, type=multi, subtopic={Basicity}, answer={A,B,C,D}, answer_src=solved, working={A is aniline: basicity aniline > chloroaniline > nitrile; carbylamine; acetanilide less reactive; Stephen reaction}, difficulty=7, diff_parts={K2 S2 I2 C1 F1}, diff_note={concepts: aniline reaction sequence, basicity | steps: protection; halogenation; diazonium chemistry; properties}]
Compound (A) has molecular formula C$_6$H$_7$N, undergoes following reaction sequence.\par
\bfig{309.3}{SA02_p2_chem_s1_q05_1.png}

all are major organic product in above reaction sequence. Select the correct statement(s) for the reaction sequence.
\begin{qoptions}
\opt{A}{Order of basic strength is A $>$ D $>$ E .}
\opt{B}{Compound (A) on reaction with CHCl$_3$, Alc. KOH results formation of \bfigi{62.0}{SA02_p2_chem_s1_q05_2.png}}
\opt{C}{Reactivity order towards aromatic electrophilic substitution is A $>$ B.}
\opt{D}{Compound (E) on reaction with SnCl$_2$, HCl followed by H$_3$O$^+$ gives \bfigi{31.4}{SA02_p2_chem_s1_q05_3.png} as a major product.}
\end{qoptions}
\end{bankq}
\begin{bankq}[id=SA05-P1-CHEM-S1-Q11, source={Score Advanced - Test 05 (15 Jul 2021)}, paper=1, qno=11, type=single, subtopic={Tests}, answer={A}, answer_src=key, difficulty=5, diff_parts={K2 S1 I2 C1 F0}, diff_note={concepts: distinguishing a free amine from its salt, tests | steps: alkali liberates the free amine so Hinsberg reagent reacts with both}]
Which of the following reagent can not be used to distinguish p-chloro aniline and anilinium hydrochloride?
\begin{qoptions}
\opt{A}{PhSO$_2$Cl + NaOH}
\opt{B}{AgNO$_3$}
\opt{C}{NaHCO$_3$}
\opt{D}{All}
\end{qoptions}
\end{bankq}
\begin{bankq}[id=SA06-P1-CHEM-S3-Q06, source={Score Advanced - Test 06 (04 Sep 2021)}, paper=1, qno=6, type=numerical, subtopic={Basicity}, answer={3}, answer_src=key, difficulty=6, diff_parts={K2 S1 I2 C1 F1}, diff_note={concepts: number of methyl iodide molecules in exhaustive methylation, basicity | steps: each methylation of the amine; Hofmann elimination}]
\bfig{121.8}{SA06_p1_chem_s3_q06_1.png}
Number of CH$_3$ -- I consumed during the following reaction. In which Hoffmann exhaustive elimination take place.
\end{bankq}
\begin{bankq}[id=SA07-P1-CHEM-S1-Q05, source={Score Advanced - Test 07 (16 Apr 2023)}, paper=1, qno=5, type=multi, subtopic={Basicity}, answer={A,B,C}, answer_src=key, difficulty=6, diff_parts={K2 S1 I2 C1 F1}, diff_note={concepts: routes to a secondary amine, basicity | steps: reductive amination of cyclopentanone with ethylamine and Gabriel-type alkylation}]
Among the following, which reaction will give N-ethylcyclopentylamine as major product?
\begin{qoptions}
\opt{A}{\bfigi{169.4}{SA07_p1_chem_s1_q05_A.png}}
\opt{B}{\bfigi{172.0}{SA07_p1_chem_s1_q05_B.png}}
\opt{C}{\bfigi{205.8}{SA07_p1_chem_s1_q05_C.png}}
\opt{D}{\bfigi{213.1}{SA07_p1_chem_s1_q05_D.png}}
\end{qoptions}
\end{bankq}

\section*{Nurture SRG papers}
\begin{bankq}[id=NSRG03-P2-CHEM-S2-Q09, source={Nurture SRG - Test 03 (04 Feb 2024)}, paper=2, qno=9, type=multi, subtopic={Basicity}, answer={A,B,C,D}, answer_src=solved, working={Both nitrogens of pyrimidine are equivalent; imidazole protonates on the pyridine type nitrogen; purine has three basic nitrogens; all three are aromatic}, difficulty=5, diff_parts={K1 S2 I2 C0 F1}, diff_note={concepts: basicity and aromaticity of heterocycles | steps: nitrogen atoms of each ring; the basic sites; aromaticity}]
\bfig{154.6}{NSRG03_p2_chem_s2_q09_1.png}
Among the following which Statements are \textbf{CORRECT}?
\begin{qoptions}
\opt{A}{Both N of pyrimidine are of same basic strength}
\opt{B}{In imidazole, protonation takes place on N---3}
\opt{C}{Purine has 3 basic Nitrogen sites}
\opt{D}{Purine, Imidazole and pyrimidine all are aromatic}
\end{qoptions}
\end{bankq}

\section*{Enthusiast Pre-Score papers}
\begin{bankq}[id=EPS01-P1-CHEM-S1-Q03, source={Enthusiast Pre-Score - Test 01 (04 Nov 2020)}, paper=1, qno=3, type=single, subtopic={Tests}, answer={C}, answer_src=key, difficulty=1, diff_parts={K1 S1 I0 C0 F0}, diff_note={concepts: isocyanide test | steps: primary amines only}]
Which of the following gives Isocyanide test with CHCl$_3$ + KOH :
\begin{qoptions}
\opt{A}{N-methylcyclohexanamine}
\opt{B}{Acetamide}
\opt{C}{Aniline}
\opt{D}{N,N-dimethylaniline}
\end{qoptions}
\end{bankq}
\begin{bankq}[id=EPS02-P1-CHEM-S3-Q04, source={Enthusiast Pre-Score - Test 02 (11 Nov 2020)}, paper=1, qno=4, type=numerical, subtopic={Tests}, answer={3.00}, answer_src=key, difficulty=7, diff_parts={K2 S2 I2 C1 F1}, diff_note={concepts: Hinsberg test, tests | steps: amine type; sulphonamide acidity; count}]
Identify compounds give white ppt. with Ph--SO$_2$Cl which is soluble in aqueous NaOH ?\par
\bfig{272.9}{EPS02_p1_chem_s3_q04_1.png}
\end{bankq}
\begin{bankq}[id=EPS02-P2-CHEM-S3-Q02, source={Enthusiast Pre-Score - Test 02 (11 Nov 2020)}, paper=2, qno=2, type=numerical, subtopic={Basicity}, answer={5.00}, answer_src=key, difficulty=7, diff_parts={K2 S2 I2 C1 F1}, diff_note={concepts: aniline and cyclohexylamine, basicity | steps: amine type; reactions of each; aromatic substitution; count}]
Compare the given characteristics of aniline and cyclohexanamine :\par
(i) Both are primary amine\par
(ii) Both can be acylated by RCOCl\par
(iii) Both reacts with CHCl$_3$/KOH\par
(iv) Both reacts with NaNO$_2$ + HCl at 0-5$^\circ$C\par
(v) Both reacts with PhSO$_2$Cl to give a compound which is soluble in KOH\par
(vi) Both gives coupling reaction with phenol\par
(vii) Both gives electrophilic substitution reaction\par
How many of the given characteristics are correct?
\end{bankq}

\section*{TOAS papers}
\begin{bankq}[id=TOAS14-P1-CHEM-S1-Q05, source={TOAS - Test 14 (07 Aug 2022)}, paper=1, qno=5, type=multi, subtopic={Tests}, answer={A,C,D}, answer_src=key, difficulty=2, diff_parts={K1 S1 I0 C0 F1}, diff_note={concepts: carbylamine reaction | steps: only primary amines give the foul-smelling isocyanide; amide is not an amine}]
Which of the following compounds will produce a compound with offensive smell when treated with CHCl$_3$/KOH?
\begin{qoptions}
\opt{A}{\bfigi{55.8}{TOAS14_p1_chem_s1_q05_A.png}}
\opt{B}{\bfigi{66.3}{TOAS14_p1_chem_s1_q05_B.png}}
\opt{C}{\bfigi{61.0}{TOAS14_p1_chem_s1_q05_C.png}}
\opt{D}{\bfigi{61.0}{TOAS14_p1_chem_s1_q05_D.png}}
\end{qoptions}
\end{bankq}
\begin{bankq}[id=TOAS15-P2-CHEM-S2-Q03, source={TOAS - Test 15 (22 Aug 2022)}, paper=2, qno=3, type=numerical, subtopic={Tests}, answer={7}, answer_src=solved, working={Primary amines give the carbylamine test: compounds 1, 3, 4, 5, 6, 8, 9 (2 is secondary, 7 is an amide): 7}, difficulty=2, diff_parts={K1 S1 I0 C0 F1}, diff_note={concepts: carbylamine test | steps: only primary amines respond; count them}]
Number of compounds show Carbylamine test, or Isolynise test
\bfig{196.4}{TOAS15_p2_chem_s2_q03_1.png}
\bfig{192.5}{TOAS15_p2_chem_s2_q03_2.png}
\bfig{177.5}{TOAS15_p2_chem_s2_q03_3.png}
\end{bankq}
\begin{bankq}[id=TOAS17-P1-CHEM-S1-Q03, source={TOAS - Test 17 (18 Sep 2022)}, paper=1, qno=3, type=multi, subtopic={Tests}, answer={A,B,C,D}, answer_src=key, difficulty=4, diff_parts={K2 S1 I1 C0 F1}, diff_note={concepts: Schmidt and Hofmann degradation, phenol tests | steps: identify each intermediate; test the product with bromine water and FeCl3}]
For the following sequence, choose the correct option(s) ?
\bfig{174.8}{TOAS17_p1_chem_s1_q03_1.png}
\begin{qoptions}
\opt{A}{(Q) will give positive Br$_2$, H$_2$O test}
\opt{B}{(Q) on reaction with neutral FeCl$_3$ gives violet colour}
\opt{C}{Major product (P) \& (S) are identical}
\opt{D}{Reaction sequence involves schmidt \& Hoffmann bromamide degradation reaction}
\end{qoptions}
\end{bankq}
\begin{bankq}[id=TOAS17-P2-CHEM-S1-Q07, source={TOAS - Test 17 (18 Sep 2022)}, paper=2, qno=7, type=single, subtopic={Basicity}, answer={D}, answer_src=key, difficulty=1, diff_parts={K1 S1 I0 C0 F0}, diff_note={concepts: Gabriel phthalimide synthesis | steps: only primary alkylamines are made; aryl halides do not react}]
Which of the following Amine can be synthesised in good yield by Gabriel phthalimide synthesis?
\begin{qoptions}
\opt{A}{Ph--NH$_2$}
\opt{B}{CH$_3$--N(CH$_3$)--CH$_3$}
\opt{C}{Ph--NH--CH$_3$}
\opt{D}{\bfigi{45.1}{TOAS17_p2_chem_s1_q07_D.png}}
\end{qoptions}
\end{bankq}
\begin{bankq}[id=TOAS18-P1-CHEM-S1-Q06, source={TOAS - Test 18 (30 Oct 2022)}, paper=1, qno=6, type=multi, subtopic={Tests}, answer={C,D}, answer_src=solved, working={P is benzyl bromide, T benzoic acid, U benzamide; W = LiAlH4 gives phenylmethanamine R (carbylamine gives the foul isocyanide); Q = AgNO2 gives the nitro compound reduced to R; V = AgCN gives benzyl isocyanide S}, difficulty=5, diff_parts={K2 S2 I1 C0 F1}, diff_note={concepts: reactions of benzyl bromide and amines, isocyanides | steps: identify each intermediate; reagents for each step; final amine and isocyanide}]
Correct option(s) for the following sequence of reactions is(are)
\bfig{230.6}{TOAS18_p1_chem_s1_q06_1.png}
\begin{qoptions}
\opt{A}{$Q$ = KNO$_2$, $W$ = LiAlH$_4$}
\opt{B}{$R$ = benzenamine, $V$ = KCN}
\opt{C}{$Q$ = AgNO$_2$, $R$ = phenylmethanamine}
\opt{D}{$W$ = LiAlH$_4$, $V$ = AgCN}
\end{qoptions}
\end{bankq}

\section*{TMAS papers}
\begin{bankq}[id=TMAS04-P2-CHEM-S1-Q03, source={TMAS - Test 04 (19 Nov 2023)}, paper=2, qno=3, type=multi, subtopic={Diazonium salts}, answer={A,D}, answer_src=key, difficulty=5, diff_parts={K1 S1 I2 C1 F1}, diff_note={concepts: diazonium reaction sequence | steps: identify each compound; properties}]
\bfig{182.5}{TMAS04_p2_chem_s1_q03_1.png}
Choose the correct statement regarding above reaction sequence
\begin{qoptions}
\opt{A}{Compound (A) is stable due to resonance}
\opt{B}{Compound (C) gives positive test of unsaturation}
\opt{C}{Compound (B) gives effervescence of $\mathrm{CO_2}$ with $\mathrm{NaHCO_3}$}
\opt{D}{Compound (A) forms red colour dye with $\beta$-naphthol}
\end{qoptions}
\end{bankq}

\section*{TNAS papers}
\begin{bankq}[id=TNAS07-P2-CHEM-S3-Q06, source={TNAS - Test 07 (20 Aug 2023)}, paper=2, qno=6, type=single, subtopic={Tests}, answer={2.00}, answer_src=key, difficulty=3, diff_parts={K1 S2 I1 C0 F0}, diff_note={concepts: Gabriel synthesis and isocyanide test | steps: which amines are primary aliphatic; which amines are primary; difference}]
Total number of amine among following which can be prepared in good yield by Gabriel phthalimide synthesis method = $X$.\par
Total number of amines giving isocyanide test = $Y$.
\bfig{296.8}{TNAS07_p2_chem_s3_q06_1.png}

Then what will value of $Y-X$ :
\end{bankq}

\section*{Star Review (Other) papers}
\begin{bankq}[id=STAR03-P1-CHEM-S2-Q08, source={Star Review (Other) - Test 03 (04 Dec 2023)}, paper=1, qno=8, type=multi, subtopic={Diazonium salts}, answer={A,B}, answer_src=key, difficulty=3, diff_parts={K1 S1 I1 C0 F1}, diff_note={concepts: formation of phenols | steps: identify each reaction; which give phenol}]
In which of the following reaction phenol obtained as product
\begin{qoptions}
\opt{A}{\bfigi{90.4}{STAR03_p1_chem_s2_q08_A.png}}
\opt{B}{\bfigi{119.5}{STAR03_p1_chem_s2_q08_B.png}}
\opt{C}{\bfigi{24.9}{STAR03_p1_chem_s2_q08_3.png}\ $\xrightarrow{\mathrm{KMnO_4}}$}
\opt{D}{\bfigi{54.6}{STAR03_p1_chem_s2_q08_D.png}\ $\xrightarrow{\text{Zn dust}}$}
\end{qoptions}
\end{bankq}
\end{document}
